\documentclass[10pt,aspectratio=169]{beamer}
\usepackage[spanish]{babel}
\usepackage[utf8]{inputenc}
\usepackage[T1]{fontenc}
\usepackage{lmodern}
\usepackage{listings}
\usepackage{tikz}
\usepackage{booktabs}
\usetikzlibrary{shapes.geometric, arrows.meta, positioning, fit, backgrounds}

% --- DEFINICIÓN MANUAL DEL LENGUAJE LUA EN LISTINGS ---
\lstdefinelanguage{Lua}{
  morekeywords={
    and,break,do,else,elseif,end,false,for,function,goto,if,in,local,
    nil,not,or,repeat,return,then,true,until,while,require,assert,import,pcall
  },
  sensitive=true,
  morecomment=[l]{--},
  morecomment=[s]{--[[}{]]},
  morestring=[b]",
  morestring=[b]'
}
\lstset{showstringspaces=false}

% --- PALETA DE COLORES ---
\definecolor{primarybg}{HTML}{0F172A}
\definecolor{slidebg}{HTML}{F8FAFC}
\definecolor{navyblue}{HTML}{1E293B}
\definecolor{accentblue}{HTML}{2563EB}
\definecolor{slate}{HTML}{64748B}

% Tonalidades del Diagrama
\definecolor{toneBlue}{HTML}{DBEAFE}
\definecolor{strokeBlue}{HTML}{2563EB}
\definecolor{toneAmber}{HTML}{FEF3C7}
\definecolor{strokeAmber}{HTML}{D97706}
\definecolor{toneMint}{HTML}{DCFCE7}
\definecolor{strokeMint}{HTML}{16A34A}
\definecolor{toneRose}{HTML}{FFE4E6}
\definecolor{strokeRose}{HTML}{E11D48}
\definecolor{toneIndigo}{HTML}{E0E7FF}
\definecolor{strokeIndigo}{HTML}{4F46E5}

% CONFIGURACIÓN DE TEMA BEAMER
\setbeamertemplate{navigation symbols}{}
\setbeamercolor{background canvas}{bg=slidebg}
\setbeamercolor{frametitle}{fg=navyblue, bg=slidebg}
\setbeamercolor{title}{fg=white}
\setbeamercolor{subtitle}{fg=toneBlue}
\setbeamercolor{structure}{fg=accentblue}
\setbeamercolor{normal text}{fg=navyblue}
\setbeamercolor{author}{fg=slate!30}

% ESTILOS DE CÓDIGO
\lstdefinestyle{luastyle}{
    language=Lua,
    basicstyle=\ttfamily\tiny,
    keywordstyle=\color{accentblue}\bfseries,
    stringstyle=\color{strokeMint},
    commentstyle=\color{strokeAmber}\itshape,
    backgroundcolor=\color{white},
    frame=single,
    rulecolor=\color{slate!30},
    breaklines=true
}

\lstdefinestyle{cstyle}{
    language=C,
    basicstyle=\ttfamily\tiny,
    keywordstyle=\color{strokeRose}\bfseries,
    stringstyle=\color{strokeMint},
    commentstyle=\color{strokeAmber}\itshape,
    backgroundcolor=\color{white},
    frame=single,
    rulecolor=\color{slate!30},
    breaklines=true
}

\lstdefinestyle{sqlstyle}{
    language=SQL,
    basicstyle=\ttfamily\tiny,
    keywordstyle=\color{strokeAmber}\bfseries,
    stringstyle=\color{strokeMint},
    backgroundcolor=\color{white},
    frame=single,
    rulecolor=\color{slate!30},
    breaklines=true
}

\lstdefinestyle{shstyle}{
    language=bash,
    basicstyle=\ttfamily\tiny,
    keywordstyle=\color{accentblue}\bfseries,
    stringstyle=\color{strokeMint},
    commentstyle=\color{slate}\itshape,
    backgroundcolor=\color{white},
    frame=single,
    rulecolor=\color{slate!30},
    breaklines=true
}

\title{\textbf{UniversalWare}}
\subtitle{Arquitectura Completa: Runtime LuaJIT, Nginx, Shell Scripts y Módulos Nativos C}
\author{José Carlos García Capistrán}
\date{}

\begin{document}

% -----------------------------------------------------------------------------
% SLIDE 1: Carátula
% -----------------------------------------------------------------------------
{
\setbeamercolor{background canvas}{bg=primarybg}
\begin{frame}
    \titlepage
\end{frame}
}

% -----------------------------------------------------------------------------
% SLIDE 2: Visión General del Sistema
% -----------------------------------------------------------------------------
\begin{frame}{Visión General del Sistema}
    \begin{columns}[T]
        \begin{column}{0.48\textwidth}
            \structure{\textbf{¿Qué es UniversalWare?}}
            \begin{itemize}
                \item Plataforma backend modular de alto rendimiento que integra orquestación en \textbf{C y C++} con runtime dinámico en \textbf{LuaJIT}.
                \item Proxy inverso de entrada en \textbf{Nginx} para manejo de archivos estáticos y limitación de tasa (\textit{rate limit}).
                \item Base de datos relacional MariaDB accedida nativamente por la extensión \texttt{cmariadb}.
            \end{itemize}
        \end{column}
        \begin{column}{0.48\textwidth}
            \structure{\textbf{Componentes Clave}}
            \begin{itemize}
                \item \textbf{Capa Web (\texttt{web/*}):} Microframework en Lua que procesa las rutas de la API REST.
                \item \textbf{Shell Tooling (\texttt{.sh}):} Scripts automáticos de compilación, resolución de dependencias y aislamiento con Podman.
                \item \textbf{Módulos \texttt{clibs}:} Extensiones C compiladas (\texttt{.so}) que otorgan capacidades de Sockets, JSON y SQL.
            \end{itemize}
        \end{column}
    \end{columns}
\end{frame}

% -----------------------------------------------------------------------------
% SLIDE 3: Arquitectura Global (Diagrama TikZ)
% -----------------------------------------------------------------------------
\begin{frame}[fragile]{Arquitectura de Componentes de UniversalWare}
\centering
\resizebox{0.92\textwidth}{!}{
\begin{tikzpicture}[
    node distance=0.45cm and 0.55cm,
    font=\sffamily\tiny,
    >=Stealth,
    box/.style={draw, rounded corners=3pt, align=center, inner sep=3.5pt, minimum height=0.55cm, fill=white},
    bBlue/.style={box, fill=toneBlue, draw=strokeBlue},
    bAmber/.style={box, fill=toneAmber, draw=strokeAmber},
    bMint/.style={box, fill=toneMint, draw=strokeMint},
    bRose/.style={box, fill=toneRose, draw=strokeRose},
    bIndigo/.style={box, fill=toneIndigo, draw=strokeIndigo},
    grp/.style={draw, dashed, inner sep=6pt, inner ysep=9pt, rounded corners=5pt, font=\sffamily\bfseries\tiny}
]

    % ==========================================
    % 1. BLOQUE DE ENTRADA (ARRIBA DERECHA)
    % ==========================================
    \node[bIndigo] (nginx) {Proxy Nginx\\{\tiny [nginx.conf:8081]}};
    \node[bIndigo, left=0.6cm of nginx] (browser) {Navegador / Cliente};

    % ==========================================
    % 2. CAPA C++ (ARRIBA IZQUIERDA)
    % ==========================================
    \node[bBlue, below=0.8cm of browser] (orchestrator) {Orquestador C++\\{\tiny [backend.cpp]}};
    \node[bBlue, right=of orchestrator] (scheduler) {Scheduler Multihilo\\{\tiny [Scheduler.h]}};
    \node[bBlue, right=of scheduler] (queue) {Cola de Trabajos\\{\tiny [Job.h]}};
    \node[bBlue, right=of queue] (worker) {Worker Ejecutor\\{\tiny [Worker.h]}};

    % ==========================================
    % 3. CAPA LUAJIT Y WEB API (CENTRO)
    % ==========================================
    \node[bAmber, below=0.8cm of orchestrator] (luajit) {Runtime LuaJIT};
    \node[bAmber, right=of luajit] (scripts) {Programas Lua\\{\tiny [program.main.lua]}};
    \node[bAmber, right=of scripts] (http_api) {API Dark Kitchen\\{\tiny [program.fetch.lua]}};
    \node[bAmber, right=of http_api] (handlers) {Handlers CRUD\\{\tiny [web/lua/fetch/]}};
    \node[bAmber, right=of handlers] (http_server) {Servidor HTTP\\{\tiny [chttp/server.c]}};

    % ==========================================
    % 4. CAPA NATIVA CLIBS Y PERSISTENCIA (ABAJO)
    % ==========================================
    \node[bMint, below=0.8cm of luajit] (mariadb_mod) {Cliente MariaDB\\{\tiny [cmariadb/main.c]}};
    \node[bMint, right=of mariadb_mod] (json_brg) {Puente JSON\\{\tiny [cjson.c]}};
    \node[bMint, right=of json_brg] (native_jobs) {Tareas Nativas\\{\tiny [cjob/main.c]}};
    \node[bRose, below=0.6cm of mariadb_mod] (mariadb) {MariaDB (Base de Datos)};

    % ==========================================
    % SUBGRAFOS / CONTENEDORES
    % ==========================================
    \begin{scope}[on background layer]
        \node[grp, draw=strokeBlue, fill=toneBlue!15, fit=(orchestrator)(worker)] (g1) {};
        \node[anchor=north west, font=\bfseries\tiny, color=strokeBlue, xshift=2pt, yshift=-1pt] at (g1.north west) {Orquestación Concurrente};

        \node[grp, draw=strokeAmber, fill=toneAmber!15, fit=(luajit)(http_server)] (g2) {};
        \node[anchor=north west, font=\bfseries\tiny, color=strokeAmber, xshift=2pt, yshift=-1pt] at (g2.north west) {Runtime Web y API};

        \node[grp, draw=strokeMint, fill=toneMint!15, fit=(mariadb_mod)(native_jobs)] (g3) {};
        \node[anchor=north west, font=\bfseries\tiny, color=strokeMint, xshift=2pt, yshift=-1pt] at (g3.north west) {Procesamiento Nativo (clibs)};

        \node[grp, draw=strokeRose, fill=toneRose!15, fit=(mariadb)] (g4) {};
        \node[anchor=north west, font=\bfseries\tiny, color=strokeRose, xshift=2pt, yshift=-1pt] at (g4.north west) {Persistencia};
    \end{scope}

    % ==========================================
    % CONEXIONES LIMPIAS Y DIRECTAS
    % ==========================================
    % Cliente -> Nginx -> API (Alineados verticalmente a la derecha)
    \draw[->] (browser) -- (nginx);
    \draw[->] (nginx) -- (http_api);

    % Flujo C++ Multihilo
    \draw[->] (orchestrator) -- (scheduler);
    \draw[->] (scheduler) -- (queue);
    \draw[->] (queue) -- (worker);

    % Worker -> LuaJIT (Conexión directa vertical C++ a Lua)
    \draw[->] (worker.south) -- ++(0,-0.15) -| (http_server.north);

    % Flujo Lua Web
    \draw[->] (luajit) -- (scripts);
    \draw[->] (scripts) -- (http_api);
    \draw[->] (http_api) -- (handlers);
    \draw[<->] (handlers) -- (http_server);

    % Handlers -> Clibs (Bajada limpia hacia la capa nativa)
    \draw[->] (handlers.south) -- (native_jobs.north);
    \draw[->] (native_jobs) -- (json_brg);
    \draw[->] (json_brg) -- (mariadb_mod);

    % Persistencia
    \draw[<->] (mariadb_mod) -- (mariadb);

\end{tikzpicture}
}
\end{frame}
% -----------------------------------------------------------------------------
% SLIDE 4: Capa de Proxy Inverso con Nginx
% -----------------------------------------------------------------------------
\begin{frame}[fragile]{Configuración del Proxy Inverso con Nginx}
    \textbf{\texttt{nginx.conf}} actúa como primera línea de defensa, sirviendo frontend estático y derivando el tráfico dinámico a LuaJIT.

    \begin{columns}[T]
        \begin{column}{0.5\textwidth}
\begin{lstlisting}[style=shstyle, caption={Reglas principales de nginx.conf}]
http {
    limit_req_zone $binary_remote_addr 
      zone=api_limit:10m rate=10r/s;

    server {
        listen 8080;
        
        # Archivos estaticos
        location / {
            root /var/www/html;
            try_files $uri$uri/ =404;
        }

        # Redireccion de API al backend
        location /api/ {
            limit_req zone=api_limit burst=20 nodelay;
            proxy_pass http://127.0.0.1:8081/;
            proxy_http_version 1.1;
            proxy_set_header Host $host;
        }
    }
}
\end{lstlisting}
        \end{column}
        \begin{column}{0.48\textwidth}
            \structure{\textbf{Ventajas del Esquema Nginx:}}
            \begin{itemize}
                \item \textbf{Seguridad y Rate Limiting:} Protege el puerto interno \texttt{8081} de ataques DDoS limitando a 10 req/s con ráfagas de hasta 20.
                \item \textbf{Aislamiento de Carga:} Nginx entrega imágenes, CSS y JS sin tocar el motor de LuaJIT.
                \item \textbf{Normalización de Cabeceras:} Agrega automáticamente IP real del cliente e información del host proxyficado.
            \end{itemize}
        \end{column}
    \end{columns}
\end{frame}

% -----------------------------------------------------------------------------
% SLIDE 5: Entorno de Automatización (.sh)
% -----------------------------------------------------------------------------
\begin{frame}[fragile]{Automatización y Tooling en Shell Scripts}
    El sistema se autoconfigura y compila mediante tres scripts principales:

    \begin{columns}[T]
        \begin{column}{0.33\textwidth}
            \structure{\textbf{1. \texttt{configurarentorno.sh}}}
            \begin{itemize}
                \item Detecta SO (Debian/Ubuntu).
                \item Compila LuaJIT v2.1 desde fuente.
                \item Instala LuaRocks y dependencias como \texttt{libssh2-1-dev}.
                \item Compila dinámicamente los módulos de \texttt{clibs/}.
            \end{itemize}
        \end{column}
        \begin{column}{0.33\textwidth}
            \structure{\textbf{2. \texttt{build.sh}}}
            \begin{itemize}
                \item Verifica si la CLI de \texttt{luajit} está presente.
                \item Compila el orquestador principal \texttt{backend.cpp} con C++17 y optimizaciones \texttt{-O3}.
            \end{itemize}
        \end{column}
        \begin{column}{0.33\textwidth}
            \structure{\textbf{3. \texttt{pods.sh}}}
            \begin{itemize}
                \item Proporciona wrappers para \textbf{Podman}.
                \item Conteneriza el entorno de desarrollo aislando permisos de usuario mediante \texttt{--userns=keep-id}.
            \end{itemize}
        \end{column}
    \end{columns}
\end{frame}

% -----------------------------------------------------------------------------
% SLIDE 6: Configuración de MariaDB CLI
% -----------------------------------------------------------------------------
\begin{frame}[fragile]{Configuración del Cliente SQL en MariaDB}
    Para establecer conexión nativa mediante \texttt{cmariadb}, el usuario cliente debe ser aprovisionado con el módulo de autenticación compatible:

    \begin{lstlisting}[style=sqlstyle, caption={Comandos ejecutados en MariaDB CLI}]
ALTER USER 'lua_client'@'localhost' IDENTIFIED VIA mysql_native_password USING PASSWORD('12345');
CREATE USER IF NOT EXISTS 'lua_client'@'127.0.0.1' IDENTIFIED VIA mysql_native_password USING PASSWORD('12345');
GRANT ALL PRIVILEGES ON *.* TO 'lua_client'@'127.0.0.1' WITH GRANT OPTION;
FLUSH PRIVILEGES;
    \end{lstlisting}

    \begin{itemize}
        \item \textbf{Conexión Loopback:} Se vincula a \texttt{127.0.0.1} para llamadas directas de la API C.
        \item \textbf{Método \texttt{mysql\_native\_password}:} Requisito clave para garantizar la compatibilidad binaria con la biblioteca del cliente C de MariaDB.
    \end{itemize}
\end{frame}

% -----------------------------------------------------------------------------
% SLIDE 7: Runtime Web en Lua (program.fetch.lua)
% -----------------------------------------------------------------------------
\begin{frame}[fragile]{Runtime Web: Enrutamiento y Sanitización}
\begin{columns}[T]
\begin{column}{0.52\textwidth}
\begin{lstlisting}[style=luastyle, caption={Ciclo principal en program.fetch.lua}]
import("cmariadb", "chttp", "cjson")

local controllers = {
    categorias = require("fetch.categorias"),
    pedidos    = require("fetch.pedidos")
}

chttp.listen("127.0.0.1", 8081)

while true do
    local peticion = chttp.accept()
    if peticion then
        local metodo = peticion.method:lower()
        local recurso = peticion.path:match("^/api/([^%?]+)")
        local controller = controllers[recurso]

        if controller and controller[metodo] then
            pcall(function()
                local db = conectar_db()
                controller[metodo](db, cjson, sanitizar, peticion, escape_sql)
                db:close()
            end)
        end
    end
end
\end{lstlisting}
\end{column}
\begin{column}{0.45\textwidth}
\structure{\textbf{Mecanismos de Control:}}
\begin{itemize}
    \item \textbf{Despacho de Controladores:} Mapea endpoints hacia scripts modulares dentro de \texttt{web/lua/fetch/}.
    \item \textbf{Aislamiento \texttt{pcall}:} Captura excepciones en las peticiones evitando la caída del daemon.
    \item \textbf{\texttt{sanitizar()}:} Función interna que convierte tipos \texttt{userdata} devueltos por C a tipos nativos de Lua interpretables por \texttt{cjson}.
\end{itemize}
\end{column}
\end{columns}
\end{frame}

% -----------------------------------------------------------------------------
% SLIDE 8: Controlador CRUD Modular (fetch/categorias.lua)
% -----------------------------------------------------------------------------
\begin{frame}[fragile]{Controladores CRUD Modulares (\texttt{web/lua/fetch/})}
\begin{lstlisting}[style=luastyle, caption={Implementación de handlers CRUD en Lua}]
local M = {}

function M.get(db, cjson, sanitizar, peticion, escape_sql)
    local res = assert(db:query("SELECT id_categoria, nombre FROM categorias ORDER BY id_categoria ASC;"))
    peticion:respond(200, cjson.encode(sanitizar(res)))
end

function M.post(db, cjson, sanitizar, peticion, escape_sql)
    local body = cjson.decode(peticion.body or "{}")
    local nombre = escape_sql(body.nombre)
    assert(db:query(string.format("INSERT INTO categorias (nombre) VALUES ('%s');", nombre)))
    peticion:respond(201, cjson.encode({ status = "ok", message = "Categoria creada" }))
end

function M.put(db, cjson, sanitizar, peticion, escape_sql)
    local body = cjson.decode(peticion.body or "{}")
    local id = tonumber(body.id_categoria)
    if not id then error("ID invalido o faltante") end
    local nombre = escape_sql(body.nombre)
    assert(db:query(string.format("UPDATE categorias SET nombre = '%s' WHERE id_categoria = %d;", nombre, id)))
    peticion:respond(200, cjson.encode({ status = "ok", message = "Categoria actualizada" }))
end

function M.delete(db, cjson, sanitizar, peticion, escape_sql)
    local id = tonumber(peticion.path:match("id_categoria=(%d+)"))
    assert(db:query(string.format("DELETE FROM categorias WHERE id_categoria = %d;", id)))
    peticion:respond(200, cjson.encode({ status = "ok", message = "Categoria eliminada" }))
end

return M
\end{lstlisting}
\end{frame}

% -----------------------------------------------------------------------------
% SLIDE 9: Catálogo de clibs Activos
% -----------------------------------------------------------------------------
\begin{frame}{Ecosistema de \texttt{clibs} en UniversalWare}
\centering
\small
\begin{tabular}{lll}
\toprule
\textbf{Módulo clib} & \textbf{Tecnología Base} & \textbf{Función Principal en el Sistema} \\
\midrule
\texttt{cmariadb} & C / MariaDB C API & Conexiones no bloqueantes, prepared statements, \texttt{SqlValue}. \\
\texttt{chttp} & C / Sockets POSIX & Servidor HTTP multihilo embebido para exponer API REST. \\
\texttt{cjson} & C / cJSON Engine & Serialización y parsing dinámico ultrarrápido de JSON. \\
\texttt{cjob} & C / Corrutinas & Despacho asíncrono de tareas a nivel de hilo nativo. \\
\bottomrule
\end{tabular}
\end{frame}

% -----------------------------------------------------------------------------
% SLIDE 10: Conclusión
% -----------------------------------------------------------------------------
\begin{frame}{Conclusiones y Arquitectura Final}
    \begin{itemize}
        \item \textbf{Arquitectura Modular en Capas:} Nginx actúa de escudo en la frontera, C++ administra la concurrencia multihilo, LuaJIT ejecuta la lógica de negocio y los \texttt{clibs} garantizan alta velocidad I/O.
        \item \textbf{Aprovisionamiento Automatizado:} El ecosistema `.sh` elimina fricciones de configuración resolviendo compilaciones nativas de dependencias.
        \item \textbf{Persistencia Directa y Segura:} Acceso veloz a MariaDB sin sobrecarga de ORMs intermedios.
    \end{itemize}
    
    \vspace{0.8cm}
    \centering
    \Large\textbf{\color{accentblue} UniversalWare: Potencia nativa, ligera y flexible.}
\end{frame}

\end{document}